% Author: Adam | Module: AI Job Interview Simulator | Project: SmartCareer AI
\section{Conclusion}
\label{sec:conclusion}

Le module livre les quatre fonctions annoncées --- génération de questions
ancrées sur un poste, évaluation par barème à quatre critères, proposition d'une
réponse améliorée structurée en STAR, historique des séances réouvrables --- avec
« Utiliser mon CV » comme branche sur le Module 1 et un choix de langue
français/anglais. Il le fait sans fine-tuning, par ingénierie de prompts et
retrieval seulement, sur un modèle de 7~Md quantifié servi localement, dans un
dépôt d'équipe dont il n'a modifié que trois fichiers de branchement.

Ce qui, dans le travail effectué, nous paraît solide parce que mesuré :

\begin{itemize}[itemsep=2pt]
  \item \textbf{La chaîne de décision de sortie tient.} 201 tests automatisés
        (200 passés, 1 sauté conditionnel, 13,15~s), et dans toutes les voies
        réelles explorées le JSON a été valide dès le premier essai : 10/10 en
        calibration directe, 4/4 re-mesurés pour ce rapport. Le seul incident
        observé (2 questions sur 3 renvoyées) a été traité par le chemin prévu ---
        conservation des questions valides, question de repli marquée
        \texttt{degraded}, réponse 201 --- sans retry et sans mensonge à
        l'interface.
  \item \textbf{La confidentialité est un invariant testé, pas une intention.}
        Suppression structurelle de \texttt{personal\_info}, effacement par
        expression régulière, neutralisation des marqueurs de gabarit en trois
        passes, bornes d'entrée, contrôle final avant départ. Le parcours
        navigateur a cherché six identifiants dans trois surfaces (payloads, page
        rendue, prompt journalisé) : 0 occurrence --- après un correctif réel,
        parce que la fuite existait (le titre de la génération CV).
  \item \textbf{Les choix de modèle et de réglages sont justifiés par des
        chiffres, y compris quand ils contredisent l'intention de départ.} La
        mesure a montré que la taille du prompt est un levier de latence
        négligeable (préfill $\approx$0,16~s pour 1800~tokens) et que
        \texttt{num\_predict} est le seul qui compte ; elle a montré aussi que
        20~couches déléguées au GPU consomment 3321~Mio sur 4096, ce qui rend la
        cohabitation avec le Module 1 impossible --- d'où une recommandation
        écrite, pas un réglage par habitude.
  \item \textbf{L'interface a été validée dans un vrai navigateur, et cela a
        changé le résultat.} Quatre des six défauts trouvés (blocage de la file
        après erreur, étiquette périmée, message serveur écrasé, seuil non
        expliqué) n'étaient pas atteignables par les tests unitaires.
\end{itemize}

Ce qui reste ouvert, et que nous écrivons sans l'adoucir : la qualité \emph{pédagogique} du score n'est pas démontrée au sens statistique --- quatre
appels sur une question, sans jeu d'or annoté, avec un biais de longueur visible
dans les données (§\ref{sec:tests-calibration}). La latence en CPU pur (10~minutes
par séance) interdit la démonstration sur une machine sans GPU. La libération de
mémoire ne fonctionne pas pendant une génération en cours. Le retrieval est
faible en français hors de son glossaire de 25~termes, et son seuil laisse passer
des hors-sujets mesurés à 0,21--0,29.

La suite logique, pour nous, n'est pas d'ajouter des fonctions mais de fermer
deux lacunes de méthode : constituer un jeu de réponses annotées par un recruteur
pour tester la validité du barème, et mesurer réellement la co-tenance des deux
modèles sur la carte de 4~Gio. Ce sont les points (c) et (d) des
perspectives (§\ref{sec:limites}) --- les deux seuls dont dépendent des
affirmations tenues aujourd'hui pour des hypothèses travaillées.
